%*******************************************************
% Acknowledgments
%*******************************************************
\pdfbookmark[1]{Acknowledgments}{acknowledgments}
\begin{flushright}{\slshape    
    I’m living for the first time. \\ \medskip
    --- Anonymous}
\end{flushright}
\bigskip
\begingroup
\let\clearpage\relax
\let\cleardoublepage\relax
\let\cleardoublepage\relax
\chapter*{Acknowledgments}

This dissertation is the end of a long path, more than eight years of study. More than once I found myself wondering how I ended up, in the first place, studying economics, and then writing a thesis that sits at the intersection of spatial economics, economic history, industrial organization, and development economics.

The reasons, of course, lie in my own life, and reflect the time I have spent with the people who influenced the path I chose to take. To them I owe a sincere thank you for what I have become and for what I have done.

The curiosity that brought me here is something I owe to my family. From the bottom of my heart, thank you to Ilaria, my sister, and to my parents. As I write these lines, it strikes me that what I ended up doing, economic geography, is precisely the combination of what my parents studied and, evidently, passed on to me: economics, from my father, and geography, from my mother.

My passion for history clearly comes from the stories in my family, from my uncles to my grandparents, who, like so many Europeans of the last century, lived through one of the most historically dense periods of the last millennium. A special thought goes to my grandmother Anna, who, instead of fairy tales when I was a child, used to tell me about Frederick Barbarossa's deeds; and to my great-uncle Paolo, who told me about the time he was a partisan during the Resistance against Nazi-fascism.

I also want to thank my PhD colleagues, with whom I spent entire days sharing ideas and opinions. A special thanks to Marc-Antoine Faure, who introduced me to spatial economics; to Sanmoy Mukherjee, with whom I discussed these topics extensively and shared conference trips. Thanks also to (in alphabetical order) Michele Badolato, Mario Bonfrisco, Manuel Delfino, Giacomo Gaggero, Diletta Migliaccio, Federico Mij, and to all the PhD students from other cohorts with whom I lived some of the best moments of my life, and who have certainly become more friends than colleagues.

Thank you as well to Laura, my partner for almost ten years, with whom I have shared and discussed my research countless times. Even though she is entirely outside economics, her wisdom and the way she sees the world have undoubtedly had one of the greatest impacts on who I am today, and on what I wrote.

A thank you also to my friends, Elisa, Federico, Giorgia, Luca, Rosa, Sabina, Stefano, Virginia, who helped shape the person I am today.

Finally, thank you to Antonella Primi, whose memory will not be forgotten.

\endgroup
